% 作品介绍：面向线性PCurve的一类通用样条面到二次曲面智能转化方案
% 第三章：2026-10-16 用户确认稿。此文件为章节片段，待整稿组装。
% 整稿需加载 amsmath、amssymb 宏包。
\section{条件假设及面向场景}

本方案面向能够在规定容差内恢复为目标解析母面的有限样条面片，支持平面、圆柱面、圆锥面、球面，以及作为扩展目标的规则圆环面。输入应具有有效的节点、次数和控制顶点数据；有理输入采用正权重，曲面内部不允许存在位置不连续的节点连接。

本方案的核心几何假设是：\textbf{输入面片的边界能够在目标解析母面对应的 NURBS 参数域中，以直线段 PCurve 的三维像表示或在容差内逼近。}设目标母面的 NURBS 参数化为
\[
\mathbf Q:\Omega_Q\subset\mathbb R^2\rightarrow\mathbb R^3,
\]
则第 $i$ 条目标边界可写为
\[
\boldsymbol{\ell}_i(t)=(1-t)\mathbf a_i+t\mathbf b_i,\qquad
\widehat{\mathbf c}_i(t)=\mathbf Q\!\left(\boldsymbol{\ell}_i(t)\right),
\quad t\in[0,1],
\]
其中，$\mathbf a_i,\mathbf b_i\in\Omega_Q$ 为参数域端点，$\boldsymbol{\ell}_i$ 为线性 PCurve，$\widehat{\mathbf c}_i$ 为其对应的三维边界曲线。输入边界与目标边界之间允许通过调整方向后的连续单调参数对应进行比较，不要求二者采用相同的参数化。

上述线性条件作用于\textbf{目标母面的 NURBS 参数域}，不要求三维边界本身为直线，也不等同于解析曲面的角度参数或其他原生参数之间存在线性关系。输入样条面的参数域与目标母面的参数域同样需要区分：源面边界位于其自身参数矩形的四个参数侧，并不能直接证明其满足目标参数域中的线性 PCurve 约束。等参边界是该约束的特殊情形；一般分支进一步处理目标参数域中的斜向直线边界。

当前一般路径以四个源参数侧为边界，恢复单一凸四边形参数区域，并处理满足退化判据的三角形情形。任意非线性 PCurve、多边界环及带孔裁剪不属于本文承诺的处理范围；球面极点、圆锥顶点及微小边界段需通过专门的退化处理与位置验收，不能因其几何尺度较小而直接忽略。

工程验证重点覆盖控制顶点总数不超过 200 的输入；当前一般路径固定测试数据采用三至五次样条表示。该表示范围属于当前验证条件，不是线性 PCurve 假设的数学必要条件。几何距离容差保持为 $10^{-5}\,\mathrm{mm}$，不通过放宽容差接受未通过验收的结果。微小分段仅在源段和目标段的连续三维直径上界均不超过 $10^{-8}\,\mathrm{mm}$ 时允许切向豁免，仍须完成位置验收，并明确记录豁免区间。

本方案以单面片的几何恢复与验收为处理边界。一般兜底采用内部采样距离检查及连续边界位置、切向认证，不将其解释为全域连续距离证明。统一入口以 $10\,\mathrm{ms}$ 为性能评价目标，两条计算路径共享 $50\,\mathrm{ms}$ 总预算；预算耗尽或实测超过上限均返回计算失败，不能据此断言不存在满足条件的解析表示。

\textbf{工程 $G^1$ 连续性保持假设：}以原相邻面片满足 $G^1$ 连续性、公共边界位置在容差内吻合为前提，将原曲面与目标解析曲面对应边界上同向单位法向的转角严格小于 $0.1^\circ$，作为保持相邻面片工程 $G^1$ 连续性的角度判据。该判据适用于法向存在的正则边界；双侧均转换时，还需复核转换后两侧法向夹角。此处为容差意义下的假设，不等同于严格 $G^1$ 证明，也不能以沿边切向认证替代。
